% ======================================================================
\section{Dust quantities seen by the chemistry and the level model}
\label{sec:dust}
% ======================================================================

MOCASSIN computes the dust absorption and scattering coefficients
$\kappa^{\mathrm{abs}}_\nu$, $\kappa^{\mathrm{sca}}_\nu$ [cm$^{-1}$] and the
asymmetry parameter $g_\nu$ of every cell from the grain species and the size
distribution of the input \citep{Ercolano2005}: per-grain cross sections from
Mie theory on optical constants \citep{BohrenHuffman1983} or from tabulated
efficiencies, weighted by each species' number abundance and the size
distribution and multiplied by the grain number density $N_d$ that follows
from the dust-to-gas mass ratio (\src{doc/\allowbreak{}dust.md}). It also solves the
radiative-equilibrium temperature of every species and size in every cell from
the local MC field. The bridge derives from these, per cell, the quantities
below. None is an input parameter: they follow the model's own grains
(\src{CLAUDE.md}, ``no hardcoded values''). Table~\ref{tab:dustfix} gives
their values on the two fixtures.

\begin{table}[H]
\centering
\caption{Dust quantities measured on the two production fixtures (run logs,
lines \texttt{[D073]}, \texttt{[D109]}, \texttt{[D195]}; values at the
calibration cell, uniform across each model).}\label{tab:dustfix}
\begin{tabular}{@{}lll@{}}
\toprule
Quantity & \texttt{hiipdr} (silicate) & \texttt{pn150k} (graphite) \\
\midrule
$\AV/N_{\mathrm H}$ [mag\,cm$^2$] & $5.18\times10^{-22}$ & $3.58\times10^{-22}$ \\
$N_{\mathrm H}/\AV$ [cm$^{-2}$\,mag$^{-1}$] & $1.93\times10^{21}$ & $2.79\times10^{21}$ \\
$\RFUV$ & 4.31 & 2.98 \\
FUV albedo, energy-weighted over 6--13.6\,eV & 0.45 & 0.37 \\
FUV asymmetry parameter $g$, scattering-weighted & 0.61 & 0.29 \\
grain area per H, silicate / carbonaceous [cm$^2$] & $1.02\times10^{-21}$ / 0 & 0 / $5.45\times10^{-22}$ \\
\bottomrule
\end{tabular}
\end{table}

\subsection{FUV-to-visual extinction ratio and \texorpdfstring{$\AV$}{AV} per H}

\rtn{dustPerCell} (\src{pdr\_mod.f90:\allowbreak{}1363--1436}) computes
\begin{gather}
  \kappa_V = \kappa^{\mathrm{abs}}_V + \kappa^{\mathrm{sca}}_V,\qquad
  \kappa_{\mathrm{FUV}} = \frac{\sum_{k\in\mathrm{FUV}} w_k\,(\kappa^{\mathrm{abs}}_k+\kappa^{\mathrm{sca}}_k)}{\sum_{k\in\mathrm{FUV}} w_k},
  \nonumber\\
  \RFUV = \frac{\kappa_{\mathrm{FUV}}}{\kappa_V},\qquad
  \frac{\AV}{N_{\mathrm H}} = \frac{1.0857\,\kappa_V}{n_{\mathrm H}}\ [\mathrm{mag\,cm^2}],
  \label{eq:rfuv}
\end{gather}
stored as \rtn{grid\%dustRfuv} and \rtn{grid\%dustAVfac}. The weight $w_k$ is the
fraction of source~1's emitted energy in bin $k$ (the increment of its
cumulative distribution), so that $\kappa_{\mathrm{FUV}}$ is the mean
extinction coefficient of the band-integrated field as the source emits it; a
mean weighted by bin width would over-weight the 912\,\AA{} end of the band,
where the extinction is steepest. The extinction includes scattering, as in
the definition of $\AV$, so $e^{-\RFUV\tau_V}$ is the transmission of the
source's FUV beam with every scattered photon counted as lost.

Two approximations limit what $\RFUV$ represents. It is weighted by the
spectrum of source~1 for every source, and by the spectrum as emitted, not as
it arrives: with depth the field is depleted first at the blue end of the band
(\dlog{D-349}), so the band-mean extinction of the local field is lower than
$\RFUV\kappa_V$ and $\RFUV$ overstates it increasingly with depth. In the 3D
build $\RFUV$ does not enter any photorate. It enters the combination of the
$\AV$ of several sources (Equation~\ref{eq:avradial}, where it cancels for a
single source), the grey extinction of the diffuse-site weights
(Section~\ref{sec:sites}), the fallback field $\chi\,e^{-\RFUV\AV/1.0857}$ of a
cell the MC field did not reach (\src{pdr\_mod.f90:\allowbreak{}3326}), and the seed.

\rtn{calibratePDRdust} (\src{pdr\_mod.f90:\allowbreak{}1439--1662}) runs once, at the
first PDR iteration. From the cell whose $\AV/N_{\mathrm H}$ is closest to the
density-weighted mean it derives the model-wide fallbacks for cells without
dust, \rtn{pdrAVfac} and \rtn{pdrRfuv}, and 3D-PDR's \rtn{UV\_fac}
$=\RFUV/1.0857$, which the build does not read (it is used only under
\kw{GRAINRECOMB}, which is not compiled; \src{pdr\_mod.f90:\allowbreak{}3308--3310}).
It reports the energy-weighted FUV albedo and the scattering-weighted
asymmetry parameter (Table~\ref{tab:dustfix}), and (\dlog{D-108}) overwrites the
extinction curve that 3D-PDR's \rtn{XLAMBDA} splines (\rtn{X\_GRID}, the
\citet{SavageMathis1979} table) with this model's $\kappa(\lambda)/\kappa_V$
at the bin nearest each tabulated wavelength.

\paragraph{3D-PDR's grain attenuation factor.} \rtn{SCATTER}
(\src{pdr/\allowbreak{}shield.F90:\allowbreak{}449--497}) is the fit of
\citet[their eqs.~1a--1d]{WagenblastHartquist1989} to the attenuation of an
FUV field incident on one side of a plane-parallel slab of absorbing and
scattering grains, which they derived with the spherical-harmonics method of
\citet{Flannery1980} for grains of albedo 0.3 and asymmetry parameter 0.8:
\begin{equation}
  f_g(\tau_\lambda) = \begin{cases}
     e^{-0.7514\,\tau_\lambda}, & \tau_\lambda<1,\\[2pt]
     \sum_{i=1}^{5} a_i\,e^{-k_i\tau_\lambda}, & \tau_\lambda\ge1,
  \end{cases}
  \qquad \tau_\lambda = \frac{\AV}{1.086}\,X(\lambda),
\end{equation}
with $(a_i)=(2.006,-1.438,0.7364,-0.5076,-0.0592)$,
$(k_i)=(0.8490,1.013,1.282,2.005,5.832)$ and $X(\lambda)=\tau_\lambda/\tau_V$
the extinction curve splined by \rtn{XLAMBDA}; the two branches agree at
$\tau_\lambda=1$ to $2\times10^{-4}$. The field thus falls as
$e^{-0.75\tau_\lambda}$ to $e^{-0.85\tau_\lambda}$ rather than $e^{-\tau_\lambda}$,
because forward scattering keeps photons in the beam. The MC transfer does
this exactly, with the model's own albedo and phase function, which differ
from 0.3 and 0.8 (Table~\ref{tab:dustfix}). In the 3D build the photorates are
evaluated at $\AV=0$ (Section~\ref{sec:rays}), where $f_g=1$, so \rtn{SCATTER}
matters only for the 1D build, the seed (Section~\ref{sec:seed}) and
diagnostics.

\subsection{Grain area per H nucleus for \texorpdfstring{\Htwo}{H2} formation}
\label{sec:grainarea}

The formation rate of \citet{CazauxTielens2002,CazauxTielens2004}
(Section~\ref{sec:formation}) is linear in the geometric grain cross section
per H nucleus of each material, for which 3D-PDR ships diffuse-ISM constants
($8.473\times10^{-22}$ and $7.908\times10^{-22}$\,cm$^2$ per H for silicate and
graphite, \src{pdr/\allowbreak{}h2\_form.F90:\allowbreak{}37--39}). \rtn{measureH2grainArea}
(\src{pdr\_mod.f90:\allowbreak{}1266--1338}) computes instead, once, the mean
geometric cross section per grain of each dust component, split by material,
\begin{equation}
  \bar\sigma_{\mathrm{mat}} = \sum_{s\in\mathrm{mat}} X_s \sum_{a=a_{\min,s}}^{a_{\max,s}} \pi a^2\, w(a)
  \quad[\mathrm{cm^2}],
\end{equation}
with $X_s$ the number abundance of species $s$ within the component and
$w(a)=n(a)\Delta a/\sum n(a)\Delta a$ the number fraction of grains in size
bin $a$, normalised over the whole size grid
(\src{ph\_mod.f90:\allowbreak{}999--1012}) and summed over the species' own size range.
This is the same weighting from which the opacity is built, so the area and
the extinction describe the same grains. Species are classified by their
label (\rtn{isCarbonGrain}, \rtn{isSilicateGrain},
\src{pdr\_mod.f90:\allowbreak{}1228--1247, 1341--1360}): carbonaceous if it contains \texttt{graph},
\texttt{grph}, \texttt{amc}, \texttt{carb}, \texttt{pah}, \texttt{sic},
\texttt{gr1} or \texttt{gr2}, silicate if it contains \texttt{sil}, \texttt{oliv},
\texttt{pyr}, \texttt{enst}, \texttt{forst}, \texttt{quartz} or \texttt{oss}; an
unrecognised label goes to the silicate channel with a warning.
\rtn{dustPerCell} then stores per cell
$\sigma^{\mathrm{sil}}_{\mathrm H}=(N_d/n_{\mathrm H})\,\bar\sigma_{\mathrm{sil}}$ and
$\sigma^{\mathrm{gra}}_{\mathrm H}=(N_d/n_{\mathrm H})\,\bar\sigma_{\mathrm{gra}}$
[cm$^2$ per H] (\rtn{grid\%dustSigSil}, \rtn{grid\%dustSigGra}); a dust-free cell
gets zero, not the upstream constant (\dlog{D-195}). On the fixtures the area
is 1.2 times the upstream silicate constant on \texttt{hiipdr} and 0.69 times
the graphite one on \texttt{pn150k}, each with no grains of the other material
(Table~\ref{tab:dustfix}); the formation rate follows linearly. The geometric
area is the right weight for the collision rate of H atoms with grains; the
efficiency per unit area is applied at one mean grain temperature
(Section~\ref{sec:Td}).

\subsection{Grain albedo in the Lyman--Werner band}

Cosmic-ray-induced photoreactions are divided by $1-\omega$
(Section~\ref{sec:cr}), with $\omega$ the grain albedo at the wavelengths of
the \Htwo{} Lyman--Werner emission that cosmic rays excite: a photon is then
scattered on average $1/(1-\omega)$ times before a grain absorbs it, and
can photodissociate or photoionise at each step. Instead of 3D-PDR's
constant $\omega=0.42$ (\src{pdr\_mod.f90:\allowbreak{}354}), each cell uses
(\rtn{grainAlbedoLW}, \src{update\_mod.f90:\allowbreak{}4658--4673}; handed over at
\src{pdr\_mod.f90:\allowbreak{}3691}, \dlog{D-280})
\begin{equation}
  \omega = \frac{\sum_{k\in\mathrm{LW}} \kappa^{\mathrm{sca}}_k\,\Delta\nu_k}
                {\sum_{k\in\mathrm{LW}} (\kappa^{\mathrm{abs}}_k+\kappa^{\mathrm{sca}}_k)\,\Delta\nu_k},
\end{equation}
an average flat in frequency over the window. MOCASSIN's own cosmic-ray-induced
ionisation of the linked neutrals uses the same $\omega$
(\src{update\_mod.f90:\allowbreak{}2916--2930}).

\subsection{Dust temperature}
\label{sec:Td}

MOCASSIN solves the radiative-equilibrium temperature $T_d(s,a)$ of every
species $s$ and size $a$ in every cell from the local MC field, by equating
the absorbed power (plus the heating by resonance lines and gas collisions,
less the photoelectric losses) to $4\pi\int C_{\mathrm{abs}}B_\nu(T_d)\,\dd\nu$; stochastic
heating is off by default (keyword \kw{quantumHeatGrain}; \dlog{L-8}). The
chemistry and the level model use one mean temperature per cell,
\begin{equation}
  T_d \equiv \rtn{Tdust(0,0,cell)} = \sum_s X_s \sum_{a=a_{\min,s}}^{a_{\max,s}} w(a)\, T_d(s,a)
\end{equation}
(\src{update\_mod.f90:\allowbreak{}4546--4555}), floored at $T_{\mathrm{CMB}}=2.725$\,K
because MOCASSIN transports no cosmic microwave background
(\src{pdr\_mod.f90:\allowbreak{}3338}; \dlog{L-30}). Three properties of this mean
matter for the molecules:
\begin{itemize}
\item it is weighted by grain \emph{number}, which favours the smallest
  grains, whereas the formation rate is a sum over grains weighted by
  \emph{area};
\item the size weights are those of the whole size grid, not renormalised to
  a species' size range, so a grid wider than the species' range enters with
  zero-temperature bins: on \texttt{test\_Problem}'s silicates (0.07--5\,\textmu m
  on a $10^{-4}$--5\,\textmu m grid) the mean came out at $(2$--$8)\times10^{-6}$\,K
  and no \Htwo{} formed anywhere (\dlog{H-035}). After \texttt{8fc1512} the user
  decided this: \dlog{D-414} renormalises each species' mean to its own bins and
  weights the species by the grains they have, which is exact (byte-identical
  outputs) where a species fills the grid;
\item one efficiency $\epsilon(T_d)$ is applied to the area of all sizes, where
  the grains have a spread of temperatures; since $\epsilon$ is non-linear,
  $\epsilon(\langle T_d\rangle)\ne\langle\epsilon(T_d)\rangle$.
\end{itemize}
The temperature is the second route by which the field acts on \Htwo{}, and
a steep one. Evaluating the efficiency of Section~\ref{sec:formation} with
3D-PDR's surface energies: on silicate $\epsilon\approx1$ for
$6\lesssim T_d\lesssim15$\,K, 0.75 at 25\,K, 0.59 at 30\,K, 0.37 at 40\,K and 0.25
at 50\,K; on graphite $\epsilon\approx1$ for $10\lesssim T_d\lesssim40$\,K, 0.91 at
50\,K and 0.43 at 100\,K. Below about 5.5\,K (silicate) and 9\,K (graphite)
$\epsilon$ collapses, because newly formed \Htwo{} no longer desorbs (the
term $F_1$ of Section~\ref{sec:formation} exceeds unity); the floor at
$T_{\mathrm{CMB}}$ does not prevent this. On \texttt{hiipdr}, whose silicate
grains sit at 26--37\,K in the PDR (\dlog{L-33}), $\dd\ln\epsilon/\dd\ln T_d\approx-1.3$
to $-1.7$: a 10\% error in $T_d$ is a 13--17\% error in the formation rate,
and the formation rate sets the depth of the H/\Htwo{} transition together
with the shielding. The grains of \texttt{pn150k} are 0.92--0.93 of Cloudy's
temperature through the PDR (\dlog{L-19}), where graphite's efficiency is flat.
